\begin{figure*}[t]
\centering
\includegraphics[width=0.96\textwidth]{fig_architecture.pdf}
\caption{Evaluation of one design: the layer profile feeds the accuracy,
memory, and cycle models, the workload simulator turns cycles and the
scheduler variables into latency and deadline misses, and the energy model
and constraint check produce the two objectives.}
\label{fig:architecture}
\end{figure*}

One design is evaluated by the pipeline of Fig.~\ref{fig:architecture}. The
network's layer profile, its MACs, weights, and output volumes, feeds
three models: a measured accuracy table, a memory model, and a cycle model.
The cycle model and the scheduler variables enter a workload simulator that
returns latency and deadline misses, and the energy model turns the
simulated activity into energy per inference. A constraint check then
decides whether the design is acceptable.

%=============================================================
% METHODOLOGY
%=============================================================

\subsection{Design Space}
\label{sec:space}
The network is a CIFAR-10 classifier of $L$ blocks, each a $3{\times}3$
convolution, batch normalization, ReLU, and $2{\times}2$ max pooling. The
blocks have $32$, $64$, $128$, $256$, and $256$ channels, each multiplied by
a width $w$, and are followed by global average pooling, a $128$-unit hidden
layer, and the ten-way classifier. Weight precision $b$ takes four values.
The scheduler parameters are the time slice $\tau$ of round-robin among
tasks of equal priority, the priority $\pi$ of the inference task relative to
the two periodic tasks, and the core clock $f$, restricted to rows of the
STM32F446RE datasheet (DS10693). One design is
$\mathbf{x}=(L,w,b,\tau,\pi,f)$. The ranges are in Table~\ref{tab:space};
$\tau$ is rounded up to whole $1$\,ms ticks.

\begin{table}[t]
\centering
\caption{Design Space}
\label{tab:space}
\begin{tabular}{|l|c|l|c|}
\hline
\rowcolor{tblhead}
\textbf{Parameter} & \textbf{Symbol} & \textbf{Values} & \textbf{Type} \\
\hline
Convolutional layers & $L$ & 2, 3, 4, 5 & 4 values \\
\hline
Width multiplier & $w$ & 0.25, 0.4375, 0.625, 0.8125, 1 & 5 values \\
\hline
Weight bit-width & $b$ & 4, 8, 16, 32 & 4 values \\
\hline
Time slice (ms) & $\tau$ & [1, 20] & range \\
\hline
Inference priority & $\pi$ & below, equal, above & 3 values \\
\hline
Clock (MHz) & $f$ & 30, 60, 120, 150, 180 & 5 values \\
\hline
\end{tabular}
\end{table}


\subsection{Workload and Energy}
\label{sec:workload}
Inference shares one core with two periodic tasks under fixed-priority
preemptive scheduling and a $1$\,kHz tick, the policy of
FreeRTOS~\cite{oliveira2023freertos}. A sensor task runs every $10$\,ms for
$90{,}000$ cycles at priority $3$, and a communication task every $50$\,ms for
$360{,}000$ cycles at priority $2$. The periods are harmonic~\cite{pavic2020periods},
the deadlines equal the periods, and the two cycle counts are a chosen
nominal case. Fifty inference requests arrive as a Poisson
process~\cite{he2020poisson} with a mean gap of $200$\,ms, each with $\pm5\%$
jitter on its cycle count. A single task serves them in arrival order, and the
priority levels \emph{below}, \emph{equal}, and \emph{above} are $1$, $2$,
and $4$. The simulator is event-driven and deterministic: every design
replays the same arrivals; designs differ only in cycles, clock, slice, and
priority.

A context switch costs $390$ cycles. That value derives from timings
measured for FreeRTOS on a Cortex-M4 (STM32F407)~\cite{ungurean2020rtos}: the
$4.13~\mu$s event-triggered switch minus the $1.82~\mu$s signalling call. It is an
upper bound, because the difference still contains the return from the
blocking call. No published value exists for the tick interrupt; its
$100$ cycles are an assumption. Sensitivity runs scale both costs by $0.5$
and $2$ and replace the task set (Section~\ref{sec:sensitivity}).

Energy per inference follows from datasheet currents,
\begin{equation}
E=\frac{V_{DD}\left(I_{\mathrm{run}}(f)\,t_{\mathrm{act}}+I_{\mathrm{sleep}}(f)\,t_{\mathrm{idle}}\right)}{N},
\label{eq:energy}
\end{equation}
with $V_{DD}=3.3$\,V and $N=50$ requests, over a horizon that ends when the
last request completes. The active time $t_{\mathrm{act}}$ covers
inference, both periodic tasks, context switches, and ticks. Run currents are
$10$, $18$, $32$, $41$, and $46$\,mA, and Sleep currents $4.43$ to
$11.36$\,mA, at $30$, $60$, $120$, $150$, and $180$\,MHz. Equation~(\ref{eq:energy}) reflects
how energy is the product of power and time~\cite{wu2021energy}, and
discrete operating points are the usual form of frequency scaling on such
parts~\cite{zidar2024energy}. The tabulated currents are typical values at
integer-mA resolution, measured with CoreMark-like code and not with network
kernels. Energy is therefore a model and was not measured on any board.

\subsection{Cycle and Memory Models Calibrated on Silicon}
\label{sec:calibration}
Cycles, flash, and RAM were measured through the STM32Cube.AI Developer
Cloud (ST Edge AI Core 4.0.1, target stm32f4, balanced optimization), which
benchmarks on real boards~\cite{berta2024cloud}. The farm lists no F446
board; the benchmarks therefore ran on an STM32F469I-DISCO at $180$\,MHz, the same
Cortex-M4F family with $2$\,MB of flash. Twelve architectures
($L=2$ to $5$, $w=0.25$, $0.625$, $1.0$) were exported in fp32 and int8, giving
$24$ benchmarks. Inference cycles are
\begin{equation}
C=c_{\mathrm{mac}}\,M+c_{\mathrm{first}}\,M_1+c_{\mathrm{w}}\,B_w+c_0,
\label{eq:cycles}
\end{equation}
with $M$ the MACs, $M_1$ those of the first convolution, $B_w$ the weight
bytes, and $c_0$ a fixed cost, fitted by non-negative least squares on
relative error separately for each precision. The int8 fit gives
$c_{\mathrm{mac}}=1.99$ cycles per MAC, $R^2=0.997$, and a leave-one-out
error of $3.3\%$ on average and $6.8\%$ at most; the fp32 fit gives
$c_{\mathrm{mac}}=6.56$, $R^2=0.997$, and $2.9\%$ and $5.9\%$. Cost per MAC
falls from $4.4$ cycles for the smallest int8 model to $2.2$ for the largest,
which the fixed and weight-proportional terms account for. Flash is $\sum$ weights${}\cdot b/8$
plus four bytes per bias plus a fitted kernel term, which matched the Cloud
report exactly for the weights. SRAM is a fitted activation buffer, a kernel
term, and a $32$\,KB reserve for the RTOS. The Cloud does not expose how it
plans the activation buffer, so that term is a regression on the summed
convolution output volume and the layer count. Leave-one-out error of the
int8 kernel flash and activation terms reaches $25\%$ and $23\%$, at most
about $8$\,KB in absolute terms; a design near the $128$\,KB limit carries
that uncertainty.

\begin{table}[t]
\centering
\caption{Calibration Quality}
\label{tab:calibration}
\footnotesize
\resizebox{\columnwidth}{!}{%
\begin{tabular}{|l|c|c|c|c|}
\hline
\rowcolor{tblhead}
\textbf{Quantity} & \textbf{n} & \textbf{$R^2$} & \textbf{LOO mean (\%)} & \textbf{LOO max (\%)} \\
\hline
Cycles, int8 & 12 & 0.997 & 3.3 & 6.8 \\
\hline
Cycles, fp32 & 12 & 0.997 & 2.9 & 5.9 \\
\hline
Kernel flash, int8 & 12 & 0.728 & 6.7 & 25.2 \\
\hline
Kernel flash, fp32 & 12 & 0.711 & 3.5 & 5.0 \\
\hline
Kernel RAM, int8 & 12 & 0.817 & 6.9 & 11.0 \\
\hline
Kernel RAM, fp32 & 12 & 0.947 & 4.2 & 5.6 \\
\hline
Activations, int8 & 12 & 0.928 & 10.5 & 22.6 \\
\hline
Activations, fp32 & 12 & 0.977 & 7.2 & 13.1 \\
\hline
\end{tabular}}

\vspace{2pt}
{\footnotesize Leave-one-architecture-out (LOO) relative error of the fits to the Cloud benchmarks.}
\end{table}


Only int8 and fp32 exist on the board path. Int16 cycles are the int8 cycles
scaled by $1.18$, the int16-to-int8 ratio measured on a Cortex-M4F for a
framework without SIMD int8~\cite{novac2021mcu}. Int4 cycles are scaled by
$1.14$, the 4-bit-weight overhead reported for CMix-NN~\cite{capotondi2020cmix}, measured
on a Cortex-M7. Both are literature-derived estimates and were not benchmarked;
int16 activations double the int8 buffer, and int4 uses the int8 buffer.

\subsection{Accuracy Table}
\label{sec:accuracy}
Accuracy is a table lookup over the $20$ architectures and four
precisions, as in tabular benchmarks of architecture
search~\cite{dong2021nats}, with no interpolation. Each architecture is
trained for $50$ epochs with Adam (learning rate $10^{-3}$, weight decay
$10^{-4}$, cosine decay), batch size $128$, random crops with padding $4$, and
horizontal flips. The training split of $45{,}000$ images and the validation
split of $5{,}000$ are fixed; the checkpoint with the best validation accuracy
is kept, and the test split is read once, at the end~\cite{lones2024pitfalls}.
Every convolution and linear layer is quantized after batch-norm folding.
Int16 and int8 use post-training quantization with per-channel symmetric
weights, per-tensor activations, and $1{,}000$ calibration images. Int4 uses
int4 weights with int8 activations, fine-tuned by quantization-aware
training~\cite{sher2023qat} for $9$ epochs at learning rate $10^{-4}$; $9$ is
the smallest count after which no later epoch of a $10$-epoch pilot, run on
two architectures, improved validation accuracy by $0.2$ points. Three
architectures were retrained with two more seeds. Table entries use seed $1$;
the other seeds measure noise, which has a standard deviation of $0.44$
points on average at fp32, $0.47$ at int8, and $1.26$ at int4. Validation
accuracy at int4 is below fp32 for all $20$ architectures, by $1.70$ points on
average and by $5.08$ at most. Two designs whose accuracies differ by less
than the noise are not distinguishable.

\begin{table}[t]
\centering
\caption{Accuracy Seed Noise (Points)}
\label{tab:accuracy}
\footnotesize
\resizebox{\columnwidth}{!}{%
\begin{tabular}{|l|c|c|c|}
\hline
\rowcolor{tblhead}
\textbf{Bit-width} & \textbf{Architectures} & \textbf{Mean std} & \textbf{Max std} \\
\hline
$b = 32$ & 3 & 0.44 & 0.72 \\
\hline
$b = 16$ & 3 & 0.45 & 0.73 \\
\hline
$b = 8$ & 3 & 0.47 & 0.78 \\
\hline
$b = 4$ & 3 & 1.26 & 2.41 \\
\hline
\end{tabular}}

\vspace{2pt}
{\footnotesize Seed std over 3 architectures trained with 3 seeds. Mean drop from 32 to 4 bits 1.70 points (range 0.58 to 5.08) over 20 architectures.}
\end{table}


\subsection{Problem Statement}
\label{sec:problem}
The search maximizes accuracy and minimizes energy per inference,
\begin{gather}
\max_{\mathbf{x}}\ \bigl(\mathrm{acc}(\mathbf{x}),\,-E(\mathbf{x})\bigr)
\ \ \text{s.t.}\ \ \mathrm{acc}\geq 0.75,\notag\\
\ell_{95}\leq 100\,\text{ms},\ \ r_{\mathrm{miss}}\leq 1\%,\notag\\
\text{flash}\leq 512\,\text{KB},\ \ \text{SRAM}\leq 128\,\text{KB},
\label{eq:problem}
\end{gather}
where $\ell_{95}$ is the $95$th-percentile inference latency over the $50$
requests and $r_{\mathrm{miss}}$ the fraction of periodic jobs that miss their
deadline, a small share of which is tolerated as in weakly hard real-time
analysis~\cite{pazzaglia2021weakly}. A design that breaks a memory limit skips the simulation and
returns infinite latency and energy; it still counts as one evaluation. Search
quality is the hypervolume of the accepted designs against a reference corner.
Corner A is $(0.70,\ 17{,}000\,\mu\mathrm{J})$ and corner B is
$(0.75,\ 23{,}000\,\mu\mathrm{J})$, both fixed before any search and both
reported, because rankings can depend on the corner~\cite{pang2024reference}.
Hypervolume is in units of accuracy fraction times $\mu$J and is expressed as
a percentage of the reference front's value.

\subsection{Search Methods and Protocol}
\label{sec:methods}
Six methods receive $100$ evaluations each, over seeds $1$ to $20$, through
one harness that owns the design space, the call to the evaluator, the
constraints, and the hypervolume trace, so that nothing but the proposal rule
differs. Flash and SRAM enter as one constraint, the larger of the two
utilizations, and the accuracy floor is a constraint for every method.
\emph{qNEHVI} fits one Gaussian process per outcome, starts from $12$
Sobol points, and maximizes the noisy log expected hypervolume improvement
with outcome constraints~\cite{low2024constrained,mamun2025mobo}, through
Ax~1.3.1 and BoTorch~0.18.1. \emph{qNParEGO} is the same with random
Chebyshev scalarization. Ax cannot form a hypervolume reference while no
evaluated design is acceptable, so both receive corner A as objective
thresholds; the other methods use no reference point. \emph{NSGA-II} is
pymoo~0.6.2~\cite{blank2020pymoo} with a population of $10$ and
constraint-domination, and \emph{MOTPE} is the multi-objective
tree-structured Parzen estimator~\cite{ozaki2022motpe} in Optuna~5.0.0 with its
constraint hook and $12$ random start trials. \emph{Random} samples the mixed
space uniformly. \emph{Sequential} is the ablation: a first qNEHVI stage of
$50$ evaluations searches $(L,w,b)$ with $\tau=10$\,ms, equal priority, and
$180$\,MHz, and a second stage of $50$ evaluations searches the scheduler
variables for a network chosen from the first stage's front.

Each adapter translates the one design space into its library's encoding.
The Ax methods use integer ranges over the index of each ordered choice,
a categorical variable for priority, and a continuous slice; this took $5$ to
$15$\,s per model fit, against $12$ to $30$\,s with categorical ordered
choices. Designs that break a memory limit have no latency or energy. The
Ax methods record only their accuracy and memory utilization, while
NSGA-II and MOTPE receive a large finite energy that ranks them last.

The reference front comes from evaluating $12{,}000$ designs, the full grid of
$L$, $w$, $b$, $f$, $\pi$, and ten log-spaced values of $\tau$. The metrics
are hypervolume~\cite{guerreiro2022hypervolume} at $25$, $50$, and $100$
evaluations, the evaluations needed to reach $90\%$ of the reference value,
and the inverted generational distance~\cite{cai2021igd} to the reference
front, with objectives scaled by the corner-A box. A Friedman test compares
the six methods; qNEHVI is compared with each other method by a one-sided
Wilcoxon signed-rank test with Holm correction~\cite{carrasco2020stats} and by Cliff's
$\delta$~\cite{ismail2026stats}. Each $p$-value is reported with the smallest
value that $20$ pairs can reach.

Hypervolume after $n$ evaluations is expressed as
\begin{equation}
\mathrm{HV}_\%(n)=100\,\frac{\mathrm{HV}\bigl(S_n\bigr)}{\mathrm{HV}\bigl(R\bigr)},
\label{eq:hv}
\end{equation}
where $S_n$ is the set of acceptable designs among the first $n$ evaluations and
$R$ the reference front; evaluations to $90\%$ is the smallest $n$ with
$\mathrm{HV}_\%(n)\geq90$, and a seed that never gets there counts as never.
Every stage, from the Cloud import to the tables, runs through one command
line entry point that stores each result by name, and the $120$ search runs
needed about $2.3$ hours of proposal and evaluation time per seed on a
fanless M4 laptop, in chunks of at most an hour.
